\section{Applications of integral transforms}

\subsection{Summing series}

Since $n^{-2}=\int_0^\infty te^{-nt}\,dt$, summing the geometric series inside the integral gives
\begin{equation}
\sum_{n=1}^{\infty}\frac{1}{n^2}
=\int_0^\infty t\sum_{n=1}^{\infty}e^{-nt}\,dt
=\int_0^\infty\frac{t}{e^t-1}\,dt
=\frac{\pi^2}{6}.
\end{equation}
The same result follows from the Parseval identity for the periodic function $f(x)=x$ on $[-\pi,\pi]$: its nonzero complex Fourier coefficients have magnitude $1/|n|$.

\subsection{Evaluating definite integrals}

When the required convergence conditions hold,
\begin{equation}
\int_p^\infty\bar f(q)\,dq\longleftrightarrow\frac{f(t)}{t}.
\end{equation}
For example,
\begin{equation}
\int_0^\infty\frac{\sin t}{t}\,dt
=\int_0^\infty\frac{1}{p^2+1}\,dp
=\frac{\pi}{2}.
\end{equation}
Applying the same idea to two cosine transforms yields, for $a,b>0$,
\begin{equation}
\int_0^\infty\frac{\cos(at)-\cos(bt)}{t}\,dt=\ln\frac{b}{a}.
\end{equation}
This kind of integral identity is useful because many impossible-looking integrals become transparent once one recognizes the correct transform pair.

\begin{note}
Use Parseval's identity for the periodic function $f(x)=x$ on $[-\pi,\pi]$:
\begin{equation}
\sum_{n=-\infty}^{\infty}|c_n|^2
=\frac{1}{2\pi}\int_{-\pi}^{\pi}x^2\,dx.
\end{equation}
For $n\neq0$,
\begin{equation}
 c_n=\frac{1}{2\pi}\int_{-\pi}^{\pi}xe^{-inx}\,dx=\frac{(-1)^n}{n}\,\mathrm{i},
\end{equation}
so $|c_n|^2=1/n^2$. Therefore,
\begin{equation}
2\sum_{n=1}^{\infty}\frac{1}{n^2}
=\frac{1}{2\pi}\int_{-\pi}^{\pi}x^2\,dx
=\frac{\pi^2}{3},
\end{equation}
which gives
\begin{equation}
\sum_{n=1}^{\infty}\frac{1}{n^2}=\frac{\pi^2}{6}.
\end{equation}
\end{note}

\subsection{Solving an initial-value problem}

Consider a mass $m$ attached to a spring of stiffness $k$:
\begin{equation}
mX''(t)+kX(t)=0,
\qquad X(0)=X_0,\quad X'(0)=0.
\end{equation}
Let $x(p)=\mathcal{L}\{X(t)\}$. The derivative rule gives
\begin{equation}
mp^2x(p)-mpX_0+kx(p)=0,
\end{equation}
and hence, with $\omega_0^2=k/m$,
\begin{equation}
x(p)=X_0\frac{p}{p^2+\omega_0^2}.
\end{equation}
Taking the inverse transform gives the familiar oscillation
\begin{equation}
X(t)=X_0\cos(\omega_0t).
\end{equation}
The example shows the practical advantage of the Laplace transform: initial values enter algebraically, and the differential equation becomes an equation for the transform.